% Lösungen Einheit 5
\einheitenkopf[][5 Sachzusammenhang]{Einheit 5 von 5 · Kurvenuntersuchung im Sachzusammenhang}

\erg{37}{a) (1) \quad b) (3) \quad c) (2) \quad d) (4) \quad e) (1)}

\erg{38}{a) Der Wasserstand steigt von 2 dm auf 18 dm, erst schnell, dann immer langsamer. \quad b) Er sinkt von 18 dm auf 10 dm, immer schneller. \quad c) Er sinkt weiter von 10 dm auf 2 dm, immer langsamer; am Ende ändert er sich kaum noch. \quad d) bei $t = 4$ und bei $t = 12$ \quad e) bei $t = 8$ (Wendepunkt) \quad f) 10 dm}

\erg{39}{a) $z'(t) = -3t^2 + 12t + 15 = -3(t - 5)(t + 1)$; $t = -1$ liegt nicht im Bereich; $z''(5) = -18 < 0$, $z(5) = 100$; Randwerte $z(0) = 0$, $z(7) = 56$: nach 5 Sekunden ist die Rakete mit 100 m am höchsten. \quad b) $n'(t) = -12t^2 + 72t$, $n'(6) = 0$, $n''(6) = -72 < 0$; $n(6) = 492$; Randwerte $n(0) = n(9) = 60$: um 15 Uhr sind 492 Besucher da, so viele wie sonst nie. \quad c) $\vartheta'(t) = 0{,}3(t - 2)(t - 6)$; Hochpunkt $(2\,|\,15{,}2)$, Tiefpunkt $(6\,|\,12)$; Randwerte $\vartheta(0) = 12$, $\vartheta(10) = 28$: höchste Temperatur 28\,°C um 18 Uhr (am Rand). \quad d) $G$ ist der Gewinn pro Tag in Euro; $G' = 0$ sucht die Stückzahl mit dem größten Gewinn, $-10$ ist keine Stückzahl; $G''(30) < 0$ zeigt, dass es ein Maximum ist; bei 30 Stück ist der Gewinn mit 170\,€ am größten ($G(0) = -100$, $G(50) = -150$). Aufgabe: „Bestimme die Stückzahl, bei der der Gewinn am größten ist, und den größten Gewinn.“}

\erg{40}{a) $w'(t) = 1{,}5t^2 - 12t + 18$, $w''(t) = 3t - 12 = 0$ für $t = 4$; $w'(4) = -6$; Randwerte $w'(0) = w'(8) = 18$: nach 4 Stunden sinkt der Wasserstand am stärksten, um 6 cm pro Stunde (bei $w(4) = 208$ cm). \quad b) $p'(t) = 30e^{-t/3}\,(1 - \tfrac{t}{3})$, $p''(t) = 30e^{-t/3}\,(\tfrac{t}{9} - \tfrac{2}{3}) = 0$ für $t = 6$; $p'(6) = -30e^{-2} \approx -4{,}06$: um 12 Uhr nimmt die Konzentration am stärksten ab, um etwa 4,06 Pollen pro m$^3$ je Stunde. \quad c) $f''(x) = -0{,}003x + 0{,}06 = 0$ für $x = 20$; $f'(20) = 0{,}6$; am Rand $f'(0) = f'(40) = 0$; größter Neigungswinkel $\tan^{-1}(0{,}6) \approx 31{,}0^\circ > 30^\circ$: Die Bedingung ist nicht erfüllt.}

\erg{41}{a) Nullstellen 0 und 4: Breite 4 LE = 20 cm; Tiefpunkt $(2\,|\,{-}1)$: Tiefe 1 LE = 5 cm \quad b) Nullstellen 0 und 6: Breite 6 LE = 30 cm $\le$ 32 cm; $m'(x) = 0{,}15x\,(x - 4)$, Tiefpunkt $(4\,|\,{-}1{,}6)$: Tiefe 1,6 LE = 8 cm $\le$ 10 cm – die Mulde passt.}

\erg{42}{a) Nils sucht $V' = 0$, dort ändert sich das Wasser gerade nicht. Am stärksten nimmt es ab, wo $V'$ am kleinsten ist: $V''(t) = 6t - 30 = 0$ für $t = 5$, $V'(5) = -12$ (Randwerte $V'(0) = 63$, $V'(8) = 15$): nach 5 Stunden, 12 m$^3$ pro Stunde. \quad b) Lea vergisst den Maßstab: 5 LE $\cdot$ 20 cm = 100 cm.}

\erg{43}{a) $V'(t)$ ist die Änderung von $V$ geteilt durch die Änderung von $t$ – Liter geteilt durch Minuten. \quad b) Steigt der Wasserstand bis 18 Uhr, ist der größte Wert am Rand des Zeitraums; dort muss die Tangente nicht waagerecht sein.}

\erg{44}{a) $N''(t) = -0{,}6t + 3 = 0$ für $t = 5$, $N'''(t) = -0{,}6 \ne 0$: am fünften Tag \quad b) $N'(t) = -0{,}3t^2 + 3t$, $N'(5) = 7{,}5$: Am fünften Tag werden etwa 7500 Downloads pro Tag erreicht, so viele wie an keinem anderen Tag. \quad c) falsch: $N'(t) = -0{,}3t\,(t - 10) > 0$ für $0 < t < 10$; die Gesamtzahl steigt weiter, nur die Downloads pro Tag werden weniger.}
